\chapter{La perte de V-JEPA~2, expliquée symbole par symbole}
\label{chap:perte-vjepa2}

\begin{quote}\itshape
Nous avons monté la machine au chapitre précédent. Nous savons qui regarde
quoi, qui devine, et qui donne le corrigé. Il reste à écrire, noir sur
blanc, le nombre que l'entraînement fait baisser. Ce nombre s'écrit en une
seule ligne de symboles, et cette ligne fait peur à beaucoup de monde. Nous
allons la découper en morceaux, et donner à chaque symbole sa ligne à lui.
Aucun ne sera laissé de côté. Et tout à la fin, nous découvrirons dans cette
ligne un tout petit détail, presque invisible, qui est le sujet de tout le
reste du livre.
\end{quote}

\begin{prerequis}
Le chapitre~\ref{chap:idee-jepa}, et lui seul. Voici ce qu'il faut en avoir
retenu. Il faut savoir qu'une vidéo est regardée deux fois, une
fois trouée et une fois complète ; que la version trouée s'appelle $x$ et la
version complète $y$ ; que l'encodeur élève $E_{\theta}$ regarde $x$, que
l'encodeur professeur $E_{\bar{\theta}}$ regarde $y$, et que le prédicteur
$P_{\varphi}$ fait le pont entre les deux. Il faut aussi se souvenir du
chapitre~\ref{chap:apprendre} : une perte est un nombre qui dit de combien
la machine se trompe. Les seules opérations employées ici seront
l'addition, la soustraction et la division.
\end{prerequis}

Au chapitre précédent, nous avons dit une phrase un peu vague, et nous
allons enfin la remplacer par quelque chose de précis. Cette phrase, la
voici : « on compare la prédiction à la cible, et l'écart entre les deux est
la perte ».

C'est une bonne phrase, mais elle ne suffit pas pour écrire un programme. Comparer,
comment ? Sur quoi, exactement ? Et pour quels morceaux de la vidéo ? Tant
que ces trois questions restent sans réponse, personne ne peut coder quoi
que ce soit.

La réponse à ces trois questions tient en une seule ligne de symboles. Cette
ligne s'appelle, dans l'article, l'équation numéro un. Nous allons la
construire ensemble, morceau par morceau, en commençant par le début.

\section{Pourquoi on écrit une formule plutôt qu'une phrase}

\begin{enclair}
Une formule n'est rien d'autre qu'une phrase, écrite en abrégé. Chaque
symbole remplace un mot ou un groupe de mots. On l'écrit ainsi pour deux
raisons : c'est plus court, et surtout c'est plus précis. Une phrase peut
souvent se comprendre de deux façons différentes. Une formule, non.
\end{enclair}

Prenons un exemple loin de notre sujet, pour bien voir l'idée. Imaginez que
je vous dise : « pour connaître le prix, on prend le nombre de pommes et on
compte deux euros chacune ». C'est clair, mais c'est long. On peut écrire la
même chose ainsi :

\[
  \text{prix} = 2 \times \text{nombre de pommes} .
\]

Rien n'a changé. On a seulement remplacé des mots par des signes. Personne
ne trouve cette ligne effrayante, et pourtant c'est déjà une formule.

Maintenant, essayons une phrase un peu plus longue : « on regarde de combien
la machine s'est trompée à chaque endroit caché, on additionne toutes ces
erreurs, et on divise par le nombre d'endroits cachés ». Cette phrase-là
tient, elle aussi, en une ligne de symboles. C'est exactement la ligne que
nous allons écrire.

\begin{intuition}
Une formule, c'est une recette de cuisine écrite en abrégé. La recette
dit : « prendre trois œufs, les battre, ajouter cent grammes de sucre ». La
formule dit la même chose avec des signes. Si vous savez lire les signes,
vous savez lire la recette. Et personne n'est né en sachant lire les signes :
cela s'apprend, un signe à la fois. C'est ce que nous allons faire.
\end{intuition}

\begin{piege}
Une formule ne cache jamais d'idée nouvelle. Elle ne fait que résumer des
idées que l'on a déjà expliquées avec des mots. Si une formule vous
paraît incompréhensible, ce n'est presque jamais parce que l'idée est
difficile : c'est parce qu'un des symboles ne vous a pas été présenté. Le
bon réflexe n'est donc pas de relire la formule dix fois, mais de repérer le
symbole que l'on ne connaît pas, et d'aller chercher ce qu'il veut dire.
\end{piege}

Voilà pourquoi ce chapitre est bâti d'une drôle de façon. Nous allons
d'abord raconter la formule en français, tranquillement. Ensuite nous
l'écrirons. Et ensuite seulement nous prendrons chaque symbole, un par un,
pour lui donner une ligne à lui tout seul.

Commençons par la version en français.

\section{La formule, d'abord racontée avec des mots}

\begin{enclair}
La machine a donné une réponse pour chaque endroit caché. Pour chacun de ces
endroits, on regarde de combien elle s'est trompée. On additionne toutes ces
erreurs. Puis on divise par le nombre d'endroits cachés, pour obtenir une
moyenne. Ce nombre-là est la perte.
\end{enclair}

Un mot, avant de commencer, sur ce que veut dire « faire une moyenne ». Vous
le faites déjà sans y penser. Trois enfants ont des billes : le premier en a
$2$, le deuxième $4$, le troisième $6$. On les met toutes en tas, ce qui en
fait $12$, puis on partage ce tas en trois parts égales. Chacun repart avec
$4$ billes. Voilà une moyenne : on met tout ensemble, puis on redistribue en
parts égales.

Pourquoi diviser, ici, plutôt que de garder la simple addition ? Parce que
sans cette division, une vidéo où l'on cache beaucoup de morceaux donnerait
toujours une erreur plus grosse qu'une vidéo où l'on en cache peu — même si
la machine travaillait aussi bien dans les deux cas. La division remet tout
le monde sur un pied d'égalité.

Ce point étant réglé, reprenons la machine du chapitre~\ref{chap:idee-jepa}
et suivons-la du doigt. Vous vous souvenez peut-être de la boucle en neuf
étapes que nous avions dressée à la fin de ce chapitre-là. Tout ce que nous
allons faire ici, c'est ouvrir la septième de ces étapes — celle qui disait
simplement « on mesure l'écart entre les réponses et les cibles » — et
regarder ce qu'elle contient vraiment. Voici donc ce qui se passe, dans
l'ordre.

D'abord, la vue trouée $x$ entre dans l'encodeur élève. Celui-ci ne voit que
les morceaux visibles, et il en ressort une version enrichie.

Ensuite, on ajoute les jetons de masque, ces enveloppes vides qui portent
l'adresse de chaque trou. Le prédicteur reçoit le tout et produit une
réponse pour chaque position.

Pendant ce temps, la vue complète $y$ est passée dans l'encodeur professeur,
qui produit les cibles, c'est-à-dire les bonnes réponses.

Enfin — et c'est le moment qui nous occupe dans ce chapitre — on prend
chaque endroit caché, on met côte à côte la réponse de la machine et la
bonne réponse, et on mesure l'écart. Puis on fait la moyenne de tous ces
écarts.

\begin{exemple}[La perte racontée sur trois trous]
Supposons, pour l'instant, qu'il n'y ait que trois endroits cachés dans
notre vidéo, et que l'on sache déjà mesurer l'écart à chacun d'eux. Disons
que ces trois écarts valent $1{,}5$, $0{,}9$ et $0{,}6$.

On les additionne : $1{,}5 + 0{,}9 + 0{,}6 = 3{,}0$. Puis on divise par
trois, puisqu'il y a trois endroits :

\[
  \frac{3{,}0}{3} = 1{,}0 .
\]

La perte vaut donc $1{,}0$. C'est tout. Le reste du chapitre ne fera
qu'expliquer comment on obtient ces trois écarts, et comment on écrit ce
calcul en abrégé.
\end{exemple}

\begin{aretenir}
Retenez la phrase, elle vous accompagnera jusqu'à la fin du livre :
\textbf{la perte est la moyenne des écarts, mesurés aux endroits cachés}.
Tout le reste n'est que de l'écriture.
\end{aretenir}

Nous avons la phrase. Écrivons-la maintenant en abrégé.

\section{La formule écrite, et son tableau de décodage}

Voici la ligne de symboles. Ne cherchez pas à la comprendre tout de suite :
regardez-la comme on regarde une photo de groupe avant de reconnaître les
visages un par un.

\begin{equation}
  \mathcal{L}_{\text{predict}}
  \;=\;
  \frac{1}{|M|}
  \sum_{i \in M}
  \Bigl\|
    \; P_{\varphi}\bigl(E_{\theta}(x),\, \Delta_y\bigr)_i
    \;-\;
    \mathrm{sg}\bigl(E_{\bar{\theta}}(y)_i\bigr) \;
  \Bigr\|_{1}
  \label{eq:perte-vjepa2}
\end{equation}

Voilà. C'est la formule complète de V-JEPA~2, celle que l'article appelle
son équation numéro un, recopiée sans rien enlever.

Elle contient exactement treize symboles différents. Pas un de plus. Et vous
en connaissez déjà la plupart, puisque nous les avons rencontrés dans les
chapitres précédents. Voici le tableau complet, sans en laisser un seul de
côté : gardez-le sous les yeux, nous allons le parcourir ligne par ligne.

\begin{center}
\footnotesize
\begin{tabular}{llp{6.1cm}}
\toprule
\textbf{Symbole} & \textbf{Se lit} & \textbf{Veut dire} \\
\midrule
$\mathcal{L}_{\text{predict}}$ & « L de prédiction » &
le nom du nombre que l'on calcule : c'est la perte. \\[3pt]

$\dfrac{1}{|M|}$ & « un sur taille de M » &
on divise par le nombre d'endroits cachés : c'est ce qui fait la
moyenne. \\[3pt]

$\displaystyle\sum_{i \in M}$ & « somme, pour $i$ dans M » &
« refais le calcul pour chaque endroit caché, et additionne tout ». \\[3pt]

$x$ & « iks » &
la vue trouée : la vidéo dont on a retiré des morceaux. \\[3pt]

$y$ & « igrec » &
la vue complète : la même vidéo, entière. \\[3pt]

$E_{\theta}$ & « E de thêta » &
l'encodeur élève, celui qui apprend. Ses boutons s'appellent $\theta$. \\[3pt]

$E_{\bar{\theta}}$ & « E de thêta barre » &
l'encodeur professeur, la copie lente de l'élève. \\[3pt]

$P_{\varphi}$ & « P de phi » &
le prédicteur, la pièce qui devine. Ses boutons s'appellent $\varphi$. \\[3pt]

$\Delta_y$ & « delta y » &
les jetons de masque : des enveloppes vides qui portent l'adresse des
trous. \\[3pt]

$\mathrm{sg}$ & « esse-gé » &
le gradient bloqué : « ne fais remonter aucune correction par ce
chemin ». \\[3pt]

$M$ & « M » &
la liste des endroits cachés. \\[3pt]

$|M|$ & « taille de M » &
combien il y a d'endroits cachés dans cette liste. \\[3pt]

$\|\,\cdot\,\|_{1}$ & « norme un » &
la façon de mesurer un écart : on additionne les différences sans se
soucier du signe. \\
\bottomrule
\end{tabular}
\end{center}

Il reste un tout petit signe qui ne figure pas dans ce tableau, parce qu'il
n'a pas de nom à lui : le minuscule $i$ posé en bas à droite de
$P_{\varphi}(\ldots)_i$ et de $E_{\bar{\theta}}(y)_i$. Nous lui donnerons
aussi son explication, un peu plus loin. Rien, vraiment rien, ne sera passé
sous silence.

Avant d'aller plus loin, arrêtons-nous une seconde sur le nom écrit à
gauche du signe égal, car il annonce déjà la suite du livre.

\begin{notation}[Le nom de la perte, $\mathcal{L}_{\text{predict}}$][nom-perte]
Nous avions vu au chapitre~\ref{chap:apprendre} que l'on note souvent une
perte par la lettre $\mathcal{L}$, d'après le mot anglais \emph{loss} qui
veut dire « perte ». Ici, cette lettre porte en plus une petite étiquette
écrite en bas : \emph{predict}, c'est-à-dire « prédiction ».

Pourquoi préciser ? Parce que cette perte-ci mesure une chose bien
particulière : ce que la machine a répondu \emph{aux endroits qu'elle devait
prédire}, c'est-à-dire aux endroits cachés. Écrire l'étiquette, c'est déjà
reconnaître qu'il pourrait en exister d'autres, mesurant autre chose.

Dans V-JEPA~2, il n'y en a pas d'autre : cette perte est la seule. Mais
gardez cette étiquette en mémoire, car au chapitre~10 nous en ajouterons une
seconde, et il faudra bien pouvoir les distinguer.
\end{notation}

\begin{exemple}[Une étiquette, c'est juste un prénom]
Dans une famille où tout le monde s'appelle Martin, on ajoute le prénom pour
savoir de qui l'on parle : Martin Paul, Martin Sophie. Le nom de famille ne
change pas, seul le prénom distingue.

C'est exactement ce qui se passe ici. Le nom de famille est $\mathcal{L}$,
et il veut dire « perte ». L'étiquette $\text{predict}$ est le prénom, et il
dit \emph{quelle} perte. Écrire
\[
  \mathcal{L}_{\text{predict}} = 0{,}5
\]
se lit donc : « la perte des endroits à prédire vaut $0{,}5$ ». Nous
calculerons ce nombre-là dans quelques pages.
\end{exemple}

Les trois dernières lignes du tableau sont les seules vraiment nouvelles :
la norme L1, l'ensemble $M$ et sa taille $|M|$. Nous leur consacrerons les
deux sections qui suivent. Mais avant, faisons un premier tour de la formule
à grande vitesse, pour voir sa forme d'ensemble.

\begin{enclair}
Lisez la formule de la droite vers la gauche, comme on démonte une machine
en commençant par l'extérieur.

Tout à droite, entre les deux grandes barres, il y a une soustraction : la
réponse de la machine moins la bonne réponse. C'est l'écart.

Les deux grandes barres autour disent comment on transforme cet écart en un
seul nombre positif.

Le grand signe $\sum$, à gauche, dit qu'on additionne ce nombre pour tous
les endroits cachés.

Et la fraction $\frac{1}{|M|}$, tout à gauche, dit qu'on divise le total par
le nombre d'endroits, pour obtenir une moyenne.

Voilà. La formule dit exactement la même chose que notre phrase en français.
\end{enclair}

La figure~\ref{fig:formule-decoupee} montre ce découpage en quatre morceaux.

\begin{schema}{La formule découpée en quatre morceaux. Les pastilles
numérotées donnent l'ordre dans lequel il faut la lire : on commence par la
droite, avec la soustraction, et l'on remonte vers la gauche jusqu'à la
moyenne. La légende du dessous reprend les quatre numéros et dit, pour
chacun, ce que ce morceau fabrique.}[formule-decoupee]

% ================== la formule, morceau par morceau ==================
% On place d'abord l'écart, qui est le plus large, puis on accroche les
% autres morceaux à sa gauche et à sa droite.
\node[font=\small, inner sep=2pt] (ecart) at (2.6,0)
     {$P_{\varphi}\bigl(E_{\theta}(x), \Delta_y\bigr)_i
       \;-\; \mathrm{sg}\bigl(E_{\bar{\theta}}(y)_i\bigr)$};
\node[font=\normalsize, inner sep=2pt, anchor=east]
      (bg) at ($(ecart.west)+(-0.18,0)$) {$\Bigl\|$};
\node[font=\normalsize, inner sep=2pt, anchor=west]
      (bd) at ($(ecart.east)+(0.18,0)$) {$\Bigr\|_{1}$};
\node[font=\normalsize, inner sep=2pt, anchor=east]
      (som) at ($(bg.west)+(-0.22,0)$) {$\displaystyle\sum_{i \in M}$};
\node[font=\normalsize, inner sep=2pt, anchor=east]
      (moy) at ($(som.west)+(-0.22,0)$) {$\dfrac{1}{|M|}$};

% ---- les fonds colorés, posés derrière ----
\begin{scope}[on background layer]
  \node[fit=(moy), inner sep=3pt, rounded corners=2pt,
        fill=mbBlue!12, draw=mbBlue!45, line width=0.5pt] (fmoy) {};
  \node[fit=(som), inner sep=3pt, rounded corners=2pt,
        fill=mbTeal!12, draw=mbTeal!45, line width=0.5pt] (fsom) {};
  \node[fit=(bg)(bd), inner sep=3pt, rounded corners=2pt,
        fill=mbAmber!10, draw=mbAmber!55, line width=0.5pt] (fbar) {};
  \node[fit=(ecart), inner sep=2pt, rounded corners=2pt,
        fill=mbViolet!10, draw=mbViolet!45, line width=0.5pt, dashed] (fec) {};
\end{scope}

% ---- les pastilles numérotées, au-dessus de chaque morceau ----
\newcommand{\mbPastille}[3]{%
  \node[circle, fill=#2, text=white, font=\sffamily\bfseries\scriptsize,
        inner sep=1.6pt] at (#1) {#3};}
\mbPastille{moy.north |- 0,0.72}{mbBlue}{4}
\mbPastille{som.north |- 0,0.72}{mbTeal}{3}
\mbPastille{bd.north |- 0,0.72}{mbAmber!85!black}{2}
\mbPastille{ecart.north |- 0,0.72}{mbViolet}{1}

% ---- le sens de lecture ----
\draw[-{Stealth[round,length=2.4mm,width=1.8mm]}, draw=mbSlate!70,
      line width=0.7pt] (5.35,1.15) -- (-1.55,1.15);
\node[font=\sffamily\scriptsize\itshape, text=mbSlate] at (1.9,1.42)
     {on lit la formule dans ce sens};

% ================== la légende ==================
\node[circle, fill=mbViolet, text=white, font=\sffamily\bfseries\scriptsize,
      inner sep=1.6pt] at (-1.55,-1.30) {1};
\node[font=\sffamily\scriptsize, text=mbInk, anchor=west] at (-1.25,-1.30)
     {\textbf{la soustraction} — ce que la machine a répondu, moins ce
      qu'il fallait répondre};

\node[circle, fill=mbAmber!85!black, text=white,
      font=\sffamily\bfseries\scriptsize, inner sep=1.6pt] at (-1.55,-1.85) {2};
\node[font=\sffamily\scriptsize, text=mbInk, anchor=west] at (-1.25,-1.85)
     {\textbf{les deux barres} — on efface les signes et on additionne :
      l'écart devient un seul nombre positif};

\node[circle, fill=mbTeal, text=white, font=\sffamily\bfseries\scriptsize,
      inner sep=1.6pt] at (-1.55,-2.40) {3};
\node[font=\sffamily\scriptsize, text=mbInk, anchor=west] at (-1.25,-2.40)
     {\textbf{le signe somme} — on recommence pour chacun des endroits
      cachés, et on additionne tout};

\node[circle, fill=mbBlue, text=white, font=\sffamily\bfseries\scriptsize,
      inner sep=1.6pt] at (-1.55,-2.95) {4};
\node[font=\sffamily\scriptsize, text=mbInk, anchor=west] at (-1.25,-2.95)
     {\textbf{la fraction} — on divise par le nombre d'endroits cachés :
      c'est la moyenne};
\end{schema}

Nous avons vu la forme d'ensemble. Approchons-nous maintenant, et regardons
de tout près la pièce qui se trouve au milieu de la formule : les deux
grandes barres.

\section{Zoom sur les deux grandes barres : la norme L1}

\begin{enclair}
Les deux grandes barres sont une machine à transformer un écart en un seul
nombre positif. Elle travaille en deux temps : on efface les signes moins, et
on additionne. C'est tout.
\end{enclair}

Avant tout le reste, il y a un petit problème à régler. Ce que le
prédicteur produit n'est pas un nombre :
c'est un \emph{jeton}, c'est-à-dire une liste de nombres, comme nous l'avons
vu au chapitre~\ref{chap:vit}. La cible n'est pas un nombre non plus :
c'est aussi une liste de nombres.

Or une perte doit être \emph{un seul nombre}. On ne peut pas descendre une
vallée dans plusieurs directions à la fois. Il faut donc une machine qui
prenne deux listes et rende un seul nombre : c'est exactement le rôle des
deux grandes barres.

Cette machine fait deux choses, l'une après l'autre. Voyons la première.

\subsection*{Premier temps : effacer les signes moins}

Quand on soustrait deux listes, on obtient une liste d'écarts. Certains sont
positifs, d'autres négatifs.

\begin{exemple}[Soustraire deux listes]
Le prédicteur a répondu $(0{,}5 \ ; \ -0{,}2 \ ; \ 1{,}0)$ et la bonne
réponse était $(0{,}4 \ ; \ 0{,}1 \ ; \ 0{,}8)$. On soustrait terme à
terme, c'est-à-dire le premier avec le premier, le deuxième avec le
deuxième, et ainsi de suite :

\[
  0{,}5 - 0{,}4 = 0{,}1 , \qquad
  -0{,}2 - 0{,}1 = -0{,}3 , \qquad
  1{,}0 - 0{,}8 = 0{,}2 .
\]

Les trois écarts valent donc $(0{,}1 \ ; \ -0{,}3 \ ; \ 0{,}2)$. Remarquez
le signe moins du deuxième : sur ce nombre-là, la machine a répondu trop
\emph{bas}.
\end{exemple}

Nous avons déjà rencontré ce problème au chapitre~\ref{chap:apprendre}, avec
la machine à bonbons : si l'on additionne les écarts tels quels, un $+5$ et
un $-5$ s'annulent, et l'on croit à tort que la machine est parfaite. Il
faut donc effacer les signes moins avant d'additionner.

Là-bas, nous avions choisi de mettre chaque écart au carré, c'est-à-dire de
le multiplier par lui-même. Ici, V-JEPA~2 fait autrement, et plus simplement
encore : il prend la \emph{valeur absolue}.

\begin{definition}[La valeur absolue][valeur-absolue]
La \emph{valeur absolue} d'un nombre, c'est ce même nombre, débarrassé de
son signe moins s'il en avait un. On l'écrit avec deux petites barres
verticales autour du nombre. Ainsi $|3| = 3$, et $|-3| = 3$ aussi. On peut
la voir comme une distance : elle dit à quelle distance de zéro se trouve le
nombre, sans dire de quel côté.
\end{definition}

\begin{exemple}[Quelques valeurs absolues]
Rien de plus simple : on regarde le nombre, et s'il porte un signe moins, on
le raye.

\begin{center}
\footnotesize
\begin{tabular}{cc}
\toprule
\textbf{Le nombre} & \textbf{Sa valeur absolue} \\
\midrule
$7$        & $7$ \\
$-7$       & $7$ \\
$0{,}3$    & $0{,}3$ \\
$-0{,}3$   & $0{,}3$ \\
$0$        & $0$ \\
\bottomrule
\end{tabular}
\end{center}

Remarquez les deux premières lignes : $7$ et $-7$ donnent la même réponse.
C'est exactement l'effet que nous cherchions, puisque nous voulions que se
tromper de sept en trop coûte autant que se tromper de sept en moins.
\end{exemple}

\begin{intuition}
Vous habitez au kilomètre zéro d'une route. Un ami habite trois kilomètres
au nord, un autre trois kilomètres au sud. Si vous demandez « à quelle
distance habites-tu ? », les deux répondent « trois kilomètres ». Aucun ne
répond « moins trois ». La valeur absolue, c'est exactement cette
réponse-là : la distance, sans la direction.
\end{intuition}

Les signes moins ont disparu. Il ne reste plus qu'un geste, et c'est le plus
simple des deux.

\subsection*{Second temps : additionner}

Une fois les signes effacés, il ne reste plus qu'à additionner tous les
écarts de la liste. On obtient enfin un seul nombre, et il est forcément
positif.

Cette machine complète — effacer les signes, puis additionner — porte un
nom.

\begin{definition}[La norme L1][norme-l1]
La \emph{norme L1} d'une liste de nombres est la somme des valeurs absolues
de ces nombres. On l'écrit en entourant la liste de deux grandes barres, et
en écrivant un petit $1$ en bas à droite : $\|\,\cdot\,\|_{1}$. On la
prononce « norme un ». Le mot \emph{norme} veut simplement dire « façon de
mesurer la taille de quelque chose ».
\end{definition}

\begin{exemple}[Calculer une norme L1, en entier]
Reprenons nos trois écarts, $(0{,}1 \ ; \ -0{,}3 \ ; \ 0{,}2)$.

On efface les signes : $(0{,}1 \ ; \ 0{,}3 \ ; \ 0{,}2)$.

On additionne : $0{,}1 + 0{,}3 + 0{,}2 = 0{,}6$.

On écrit donc :
\[
  \bigl\| (0{,}1 \ ; \ -0{,}3 \ ; \ 0{,}2) \bigr\|_{1} = 0{,}6 .
\]

Trois nombres sont devenus un seul, et ce seul nombre dit de combien la
machine s'est trompée sur ce jeton-là. Retenez ce $0{,}6$ : il nous servira
tout à l'heure.
\end{exemple}

La figure~\ref{fig:norme-l1} montre les deux temps de l'opération.

\begin{schema}{La norme L1, en deux temps. À gauche, les trois écarts tels
qu'ils sortent de la soustraction : la longueur de chaque barre donne sa
taille, et celle du milieu est rouge parce qu'elle est négative. Au milieu,
les signes ont été effacés : les trois barres sont devenues des longueurs,
toutes de la même couleur. À droite, on les met bout à bout, et la longueur
obtenue est la norme L1.}[norme-l1]

% Échelle : 1 unité de valeur = 4 cm. Ainsi 0,1 -> 0,4 et 0,3 -> 1,2.
% Les trois barres sont empilées, ce qui permet de comparer leurs longueurs
% d'un coup d'œil.

% ================== 1. les écarts, avec leurs signes ==================
\node[font=\sffamily\scriptsize\bfseries, text=mbSlate] at (0.75,1.55)
     {les écarts};
\node[font=\small, text=mbInk, anchor=east] at (-0.15,1.0)  {$0{,}1$};
\node[font=\small, text=mbRose!85!black, anchor=east] at (-0.15,0.5) {$-0{,}3$};
\node[font=\small, text=mbInk, anchor=east] at (-0.15,0.0)  {$0{,}2$};

\fill[mbAzure!30, draw=mbAzure!75, line width=0.6pt]
      (0,0.85) rectangle ++(0.4,0.3);
\fill[mbRose!30, draw=mbRose!80, line width=0.6pt]
      (0,0.35) rectangle ++(1.2,0.3);
\fill[mbAzure!30, draw=mbAzure!75, line width=0.6pt]
      (0,-0.15) rectangle ++(0.8,0.3);

\node[font=\sffamily\scriptsize, text=mbRose!85!black, anchor=north west]
      at (0,-0.42) {celui-ci est négatif};

% ================== la première flèche ==================
\draw[fleche accent] (2.15,0.5) -- (3.55,0.5);
\node[font=\sffamily\scriptsize, text=mbAmber!80!black, align=center]
      at (2.85,0.95) {on efface\\les signes};

% ================== 2. les longueurs, toutes positives ==================
\begin{scope}[shift={(4.35cm,0cm)}]
  \node[font=\sffamily\scriptsize\bfseries, text=mbSlate] at (0.75,1.55)
       {des longueurs};
  \node[font=\small, text=mbInk, anchor=east] at (-0.15,1.0) {$0{,}1$};
  \node[font=\small, text=mbInk, anchor=east] at (-0.15,0.5) {$0{,}3$};
  \node[font=\small, text=mbInk, anchor=east] at (-0.15,0.0) {$0{,}2$};

  \fill[mbAzure!30, draw=mbAzure!75, line width=0.6pt]
        (0,0.85) rectangle ++(0.4,0.3);
  \fill[mbAzure!30, draw=mbAzure!75, line width=0.6pt]
        (0,0.35) rectangle ++(1.2,0.3);
  \fill[mbAzure!30, draw=mbAzure!75, line width=0.6pt]
        (0,-0.15) rectangle ++(0.8,0.3);

  \node[font=\sffamily\scriptsize, text=mbSlate, anchor=north west]
        at (0,-0.42) {plus aucun signe moins};
\end{scope}

% ================== la seconde flèche ==================
\draw[fleche accent] (6.50,0.5) -- (7.90,0.5);
\node[font=\sffamily\scriptsize, text=mbAmber!80!black, align=center]
      at (7.20,0.95) {on les met\\bout à bout};

% ================== 3. la somme ==================
\begin{scope}[shift={(8.70cm,0cm)}]
  \node[font=\sffamily\scriptsize\bfseries, text=mbSlate] at (1.2,1.55)
       {une seule longueur};

  \fill[mbAzure!30, draw=mbAzure!75, line width=0.6pt]
        (0,0.35) rectangle ++(0.4,0.3);
  \fill[mbAzure!45, draw=mbAzure!75, line width=0.6pt]
        (0.4,0.35) rectangle ++(1.2,0.3);
  \fill[mbAzure!30, draw=mbAzure!75, line width=0.6pt]
        (1.6,0.35) rectangle ++(0.8,0.3);

  \draw[accolade bas] (0,0.28) -- (2.4,0.28);
  \node[font=\sffamily\small\bfseries, text=mbAzure!70!black] at (1.2,-0.28)
       {en tout : $0{,}6$};
  \node[font=\sffamily\scriptsize, text=mbSlate] at (1.2,-0.75)
       {c'est la norme L1};
\end{scope}
\end{schema}

\begin{plusloin}
Pourquoi la valeur absolue plutôt que le carré du
chapitre~\ref{chap:apprendre} ? Parce que les deux ne punissent pas de la
même façon.

Prenons deux écarts, un tout petit et un gros : $0{,}1$ et $1{,}0$.

Avec la norme L1, le total vaut $0{,}1 + 1{,}0 = 1{,}1$. Le gros écart
compte pour $1{,}0 / 1{,}1 \approx 91\,\%$ du total.

Avec les carrés, le total vaut $0{,}01 + 1{,}00 = 1{,}01$. Le gros écart
compte cette fois pour $1{,}00 / 1{,}01 \approx 99\,\%$ du total, et le petit
pour à peine $1\,\%$.

Le carré écrase donc les petits écarts et ne s'occupe presque plus que des
gros. Or, dans une vidéo, les gros écarts arrivent souvent aux endroits
imprévisibles : une porte qui claque, quelqu'un qui entre dans le champ.
Avec le carré, la machine passerait son temps sur ces surprises qu'elle ne
pouvait pas deviner. Avec la valeur absolue, elle continue d'écouter les
petits écarts, ceux qui viennent des endroits qu'elle aurait vraiment pu
mieux comprendre.
\end{plusloin}

Voilà les deux grandes barres expliquées. Il nous reste la lettre $M$ et sa
taille $|M|$ : c'est encore plus simple, et c'est le sujet de la section qui
vient.

\section{Zoom sur la lettre M : la liste des endroits cachés}

\begin{enclair}
La lettre $M$ est la liste des adresses de tous les trous. Ce n'est pas un
nombre, c'est une liste. Et $|M|$, avec les deux barres, est la longueur de
cette liste, c'est-à-dire combien il y a de trous.
\end{enclair}

Reprenons notre clip de la phase principale, celui qui nous accompagne
depuis le chapitre~\ref{chap:vit}. Il contient $2\,048$ jetons. On leur a
donné un numéro, de $1$ jusqu'à $2\,048$, pour pouvoir en parler.

Maintenant, on tire un masque au hasard. Certains jetons sont recouverts,
d'autres restent visibles. Il faut bien noter quelque part lesquels sont
recouverts. C'est ce que fait la lettre $M$.

\begin{definition}[Un ensemble][ensemble]
Un \emph{ensemble} est simplement une liste d'objets. Chaque objet n'y
figure qu'une seule fois, et l'ordre dans lequel on les écrit n'a aucune
importance. On note un ensemble entre accolades. Par exemple,
$\{3 \ ; \ 17 \ ; \ 42\}$ est l'ensemble qui contient les trois nombres $3$,
$17$ et $42$. Dire qu'un objet « appartient » à l'ensemble, c'est dire qu'il
est dans la liste, et on l'écrit avec le petit signe $\in$, qui se lit
« appartient à ».
\end{definition}

\begin{definition}[L'ensemble des endroits cachés, noté $M$][ensemble-m]
On note $M$ l'ensemble des numéros des jetons masqués. La lettre $M$ vient
de l'anglais \emph{masked}, qui veut dire « masqué ». Écrire $i \in M$ se lit
« le jeton numéro $i$ fait partie des jetons masqués ».
\end{definition}

\begin{definition}[La taille d'un ensemble, notée $|M|$][taille-ensemble]
Les deux barres autour d'un ensemble donnent le nombre d'objets qu'il
contient. Ainsi, si $M = \{3 \ ; \ 17 \ ; \ 42\}$, alors $|M| = 3$.
Attention : ce sont les mêmes barres que celles de la valeur absolue, mais
elles ne font pas la même chose. Autour d'un nombre, elles effacent le
signe ; autour d'un ensemble, elles comptent combien il y a d'objets dedans.
C'est le contexte
qui dit laquelle des deux on emploie.
\end{definition}

\begin{exemple}[Un masque minuscule, écrit en entier]
Prenons une toute petite vidéo qui ne contiendrait que huit jetons,
numérotés de $1$ à $8$. On tire un masque, et il recouvre les jetons
numéro $3$, $4$, $5$ et $7$. On écrit alors :

\[
  M = \{3 \ ; \ 4 \ ; \ 5 \ ; \ 7\} , \qquad |M| = 4 .
\]

Les quatre jetons restants, ceux numéro $1$, $2$, $6$ et $8$, forment le
contexte. Nous leur donnerons leur propre lettre au chapitre~10 ; pour
l'instant, ils n'ont pas de nom dans la formule, et c'est justement là que
tout se joue.
\end{exemple}

\begin{exemple}[Le masque de notre clip habituel]
Reprenons les chiffres du chapitre~\ref{chap:idee-jepa}. Notre clip contient
$2\,048$ jetons, et le masque en recouvre $1\,200$. On a donc :

\[
  |M| = 1\,200 .
\]

L'ensemble $M$ contient donc mille deux cents numéros. On ne les écrit
évidemment jamais à la main : c'est le programme qui les garde en mémoire.
Les $848$ jetons restants sont le contexte.
\end{exemple}

La figure~\ref{fig:ensemble-m} met cette petite vidéo sous vos yeux, pour
que vous voyiez d'un coup ce que la lettre $M$ contient et ce qu'elle laisse
de côté.

\begin{schema}{Ce que contient la lettre $M$. Notre petite vidéo n'a que
huit jetons, numérotés de $1$ à $8$. Quatre d'entre eux sont cachés : ce sont
les gris, et leurs quatre numéros forment l'ensemble $M$. Les quatre autres,
les bleus, restent visibles : ce sont les jetons de contexte, et aucun de
leurs numéros n'entre dans $M$.}[ensemble-m]

% ---- la file des huit jetons ----
\foreach \k in {1,...,8}{
  \node[font=\sffamily\scriptsize, text=mbSlate] at ({(\k-1)*1.15+0.4},1.05)
       {\k};
}
\foreach \k in {1,2,6,8}{
  \fill[mbAzure!25, draw=mbAzure!70, line width=0.7pt]
        ({(\k-1)*1.15},0) rectangle ++(0.8,0.8);
}
\foreach \k in {3,4,5,7}{
  \fill[mbSlate!35, draw=mbSlate!80, line width=0.7pt, dashed]
        ({(\k-1)*1.15},0) rectangle ++(0.8,0.8);
}

\node[font=\sffamily\scriptsize, text=mbSlate, anchor=east] at (-0.25,0.4)
     {les huit jetons};

% ---- on descend les numeros gris dans l'ensemble ----
\foreach \k in {3,4,5,7}{
  \draw[trait annotation] ({(\k-1)*1.15+0.4},-0.12) -- (4.0,-0.95);
}
\node[font=\small, text=mbInk] at (4.0,-1.35)
     {$M = \{\,3 \ ; \ 4 \ ; \ 5 \ ; \ 7\,\}$};
\node[font=\small, text=mbInk] at (4.0,-2.0)
     {$|M| = 4$};
\node[annotation, anchor=west, align=left] at (6.4,-1.35)
     {les numéros\\des cases cachées};
\node[annotation, anchor=west, align=left] at (6.4,-2.0)
     {combien il y en a};

\node[font=\sffamily\scriptsize, text=mbAzure!80!black, anchor=west]
     at (-0.25,-2.75)
     {les jetons $1$, $2$, $6$ et $8$ restent dehors : ils sont le contexte};
\end{schema}

Il nous reste à comprendre ce que fait le grand signe $\sum$, celui qui est
posé juste devant $M$ dans la formule.

\begin{notation}[Le signe somme, revu][somme-revue]
Nous avions déjà croisé ce signe au chapitre~\ref{chap:apprendre}. Écrire
$\sum_{i \in M}$ veut dire : « fais prendre à $i$ chacune des valeurs de
l'ensemble $M$, calcule à chaque fois ce qui suit, et additionne tous les
résultats ».

Avec notre masque minuscule $M = \{3 \ ; \ 4 \ ; \ 5 \ ; \ 7\}$, cela veut
donc dire : calcule pour le jeton~$3$, puis pour le jeton~$4$, puis pour le
jeton~$5$, puis pour le jeton~$7$, et additionne les quatre résultats.
Et surtout : \textbf{ne calcule rien du tout pour les jetons $1$, $2$, $6$
et $8$}. Ils ne sont pas dans la liste, donc on ne les regarde pas.
\end{notation}

Cette dernière phrase, en gras, a l'air d'un détail d'écriture. Retenez-la :
c'est elle qui va tout faire basculer dans quelques pages.

\begin{notation}[Le petit $i$ en bas à droite][indice-i]
Vous avez sans doute remarqué le petit $i$ posé en bas à droite, dans
$P_{\varphi}(\ldots)_i$ et dans $E_{\bar{\theta}}(y)_i$. Il veut dire
« prends seulement le jeton numéro $i$ dans tout ce qui sort ».

Souvenez-vous du chapitre~\ref{chap:vit} : l'encodeur ne rend pas un seul
paquet de nombres, mais un paquet \emph{par jeton}. Le prédicteur non plus.
Le petit $i$ est donc l'instruction qui dit : « dans ces deux mille
quarante-huit paquets, va me chercher celui qui porte le numéro $i$ ».
\end{notation}

\begin{exemple}[Le petit $i$ au travail]
C'est le même geste que lorsque vous cherchez un appartement dans un
immeuble. Le prédicteur, c'est l'immeuble entier : il contient $2\,048$
appartements, un par jeton. L'écriture
$P_{\varphi}\bigl(E_{\theta}(x), \Delta_y\bigr)_{17}$ veut simplement dire :
« va sonner à la porte numéro $17$ ».

Si l'on remplace $17$ par $500$, on obtient un autre paquet de nombres,
celui d'un tout autre endroit de la vidéo. Le calcul qui a fabriqué
l'immeuble est le même ; seul le numéro de la porte a changé.
\end{exemple}

Nous avons maintenant tous les symboles. Mettons-les au travail sur de vrais
nombres.

\section{Le calcul complet, du début à la fin}
\label{sec:calcul-complet}

\begin{enclair}
Nous allons calculer une perte entière, sans rien sauter. Pour que tout
tienne sur une ligne, nos jetons n'auront que trois nombres au lieu de mille,
et notre masque ne recouvrira que deux endroits au lieu de mille deux cents.
L'idée est exactement la même ; seules les tailles changent.
\end{enclair}

Voici la situation. Notre masque recouvre deux jetons, disons ceux qui
portent les numéros $3$ et $7$. On écrit donc :

\[
  M = \{3 \ ; \ 7\} , \qquad |M| = 2 .
\]

Nous allons traiter ces deux endroits l'un après l'autre, puis rassembler
les deux résultats. Commençons par le premier.

\subsection*{Étape 1 : le jeton numéro 3}

Le prédicteur a répondu, pour cet endroit :
\[
  P_{\varphi}\bigl(E_{\theta}(x), \Delta_y\bigr)_3
  = (0{,}5 \ ; \ -0{,}2 \ ; \ 1{,}0) .
\]

Et le professeur, qui voyait la vidéo complète, avait produit comme cible :
\[
  \mathrm{sg}\bigl(E_{\bar{\theta}}(y)_3\bigr)
  = (0{,}4 \ ; \ 0{,}1 \ ; \ 0{,}8) .
\]

On soustrait terme à terme :
\[
  (0{,}5 - 0{,}4 \ ; \ -0{,}2 - 0{,}1 \ ; \ 1{,}0 - 0{,}8)
  = (0{,}1 \ ; \ -0{,}3 \ ; \ 0{,}2) .
\]

On efface les signes et on additionne :
\[
  0{,}1 + 0{,}3 + 0{,}2 = 0{,}6 .
\]

L'erreur sur le jeton numéro $3$ vaut donc $\bm{0{,}6}$. Un endroit caché
est réglé ; passons au second, en refaisant exactement les mêmes gestes.

\subsection*{Étape 2 : le jeton numéro 7}

Cette fois, le prédicteur a répondu, et le professeur avait produit sa
cible :
\[
  P_{\varphi}\bigl(E_{\theta}(x), \Delta_y\bigr)_7 = (0{,}0 \ ; \ 0{,}0 \ ; \ 0{,}0) ,
  \qquad
  \mathrm{sg}\bigl(E_{\bar{\theta}}(y)_7\bigr) = (0{,}1 \ ; \ -0{,}2 \ ; \ 0{,}1) .
\]

On soustrait :
\[
  (0{,}0 - 0{,}1 \ ; \ 0{,}0 + 0{,}2 \ ; \ 0{,}0 - 0{,}1)
  = (-0{,}1 \ ; \ 0{,}2 \ ; \ -0{,}1) .
\]

On efface les signes et on additionne :
\[
  0{,}1 + 0{,}2 + 0{,}1 = 0{,}4 .
\]

L'erreur sur le jeton numéro $7$ vaut donc $\bm{0{,}4}$. Nous avons nos deux
erreurs : il ne reste plus qu'à les rassembler en un seul nombre.

\subsection*{Étape 3 : la moyenne}

Le signe $\sum$ nous dit d'additionner les deux résultats :
\[
  0{,}6 + 0{,}4 = 1{,}0 .
\]

Et la fraction $\frac{1}{|M|}$ nous dit de diviser par le nombre d'endroits
cachés, c'est-à-dire par deux :
\[
  \mathcal{L}_{\text{predict}} = \frac{1{,}0}{2} = \bm{0{,}5} .
\]

\begin{exemple}[Le calcul rassemblé dans un tableau]
Voici les mêmes calculs, rangés d'un coup d'œil :

\begin{center}
\footnotesize
\begin{tabular}{cccc}
\toprule
\textbf{Jeton} & \textbf{Réponse} & \textbf{Cible} & \textbf{Écart, norme L1} \\
\midrule
$3$ & $(0{,}5 \ ; \ -0{,}2 \ ; \ 1{,}0)$ & $(0{,}4 \ ; \ 0{,}1 \ ; \ 0{,}8)$ & $0{,}6$ \\
$7$ & $(0{,}0 \ ; \ 0{,}0 \ ; \ 0{,}0)$ & $(0{,}1 \ ; \ -0{,}2 \ ; \ 0{,}1)$ & $0{,}4$ \\
\midrule
\multicolumn{3}{r}{\textbf{somme}} & $1{,}0$ \\
\multicolumn{3}{r}{\textbf{divisée par } $|M| = 2$} & $\bm{0{,}5}$ \\
\bottomrule
\end{tabular}
\end{center}
\end{exemple}

\begin{aretenir}
Vous venez de calculer, à la main et en entier, la perte de V-JEPA~2. Dans
l'article, les jetons contiennent mille nombres au lieu de trois, et le
masque couvre mille deux cents endroits au lieu de deux. Mais le calcul ne
fait strictement rien d'autre que ce que vous venez de faire. Si vous savez
refaire ce tableau, vous savez lire l'équation~1 de l'article.
\end{aretenir}

Nous avons fini de lire la formule. Il est temps de regarder ce qu'elle ne
dit pas.

\section{Le détail qui change tout : où la perte regarde-t-elle ?}

\begin{enclair}
Le prédicteur produit une réponse pour \emph{tous} les jetons, y compris
ceux qui étaient déjà visibles. Mais la perte, elle, ne regarde que les
jetons masqués. Toutes les autres réponses sont jetées à la poubelle, sans
être notées. Ce détail est le sujet de tout le reste de ce livre.
\end{enclair}

Revenons une dernière fois sur le signe $\sum_{i \in M}$. Nous avons dit
qu'il fait prendre à $i$ toutes les valeurs de l'ensemble $M$. Posons-nous
donc la question qui fâche : \emph{que se passe-t-il pour les jetons qui ne
sont pas dans $M$ ?}

La réponse est brutale : rien. Absolument rien. Ils ne sont pas dans la
liste, donc on ne les calcule pas, donc ils ne comptent pas dans la perte.

Or, rappelez-vous ce que nous avions compté au chapitre~\ref{chap:idee-jepa}.
Le prédicteur reçoit les jetons de contexte \emph{et} les jetons de masque,
et il ressort une réponse pour chaque jeton reçu. Il travaille donc pour
tout le monde. Mais on ne le note que sur une partie de sa copie.

\begin{exemple}[Combien de réponses partent à la poubelle]
Reprenons nos chiffres habituels.

\begin{center}
\footnotesize
\begin{tabular}{lr}
\toprule
\textbf{Ce que l'on compte} & \textbf{Combien} \\
\midrule
Réponses produites par le prédicteur & $2\,048$ \\
Réponses regardées par la perte (les masquées) & $1\,200$ \\
Réponses jetées sans être regardées (le contexte) & $848$ \\
\bottomrule
\end{tabular}
\end{center}

La part jetée vaut
\[
  \frac{848}{2\,048} \approx 0{,}414 ,
\]
c'est-à-dire un peu plus de \textbf{quarante et un pour cent}. Plus de
quatre réponses sur dix sont donc fabriquées, puis jetées sans que personne
ne les regarde.
\end{exemple}

La figure~\ref{fig:perte-regarde} est le dessin le plus important de cette
partie du livre. Regardez-le longuement, et surtout retenez ses couleurs :
nous le retrouverons au chapitre~10, presque identique, avec un seul
changement — et ce changement sera toute la nouveauté de V-JEPA~2.1.

\begin{schema}{Où la perte de V-JEPA~2 regarde-t-elle. Chaque petit carré
est un jeton, et le prédicteur a produit une réponse pour les quarante. Mais
l'accolade de droite ne ramasse que les vingt-quatre carrés gris, ceux qui
étaient cachés : eux seuls entrent dans la perte. Les seize carrés bleus
étaient visibles, et leurs réponses partent à la poubelle sans avoir été
regardées.}[perte-regarde]

% ================== la grille : 8 colonnes x 5 lignes ==================
% cellule de 0,70 ; la grille occupe donc 5,60 de large et 3,50 de haut.
\foreach \c in {0,...,7}{
  \foreach \r in {0,...,4}{
    \fill[mbAzure!28, draw=mbAzure!70, line width=0.5pt]
          (\c*0.70, -\r*0.70) rectangle ++(0.66,-0.66);
  }
}
% ---- le bloc masqué : colonnes 1 à 6, lignes 1 à 4 (24 carrés) ----
\foreach \c in {1,...,6}{
  \foreach \r in {1,...,4}{
    \fill[mbSlate!38, draw=mbSlate!75, line width=0.5pt]
          (\c*0.70, -\r*0.70) rectangle ++(0.66,-0.66);
  }
}

% ================== l'accolade : elle ne prend QUE le gris ==================
% le bloc gris va de y = -0,70 à y = -3,50 ; la grille s'arrête à x = 5,26.
\draw[accolade, line width=0.9pt, draw=mbRose!85!black]
      (5.95,-3.50) -- (5.95,-0.70);
\node[font=\sffamily\small\bfseries, text=mbRose!85!black, anchor=west,
      align=left, text width=4.6cm] at (6.50,-1.60)
     {les 24 réponses\\que la perte regarde};
\node[font=\sffamily\scriptsize, text=mbSlate, anchor=west,
      align=left, text width=4.6cm] at (6.50,-2.55)
     {ce sont elles, et elles seules,\\que désigne l'ensemble $M$};

% ================== l'accolade du bas : les 40 réponses ==================
\draw[accolade bas] (0,-3.72) -- (5.56,-3.72);
\node[font=\sffamily\small\bfseries, text=mbSlate] at (2.78,-4.32)
     {40 réponses produites par le prédicteur};

% ================== la légende, sous le dessin ==================
\fill[mbSlate!38, draw=mbSlate!75, line width=0.5pt]
      (0,-5.15) rectangle ++(0.42,-0.42);
\node[font=\sffamily\scriptsize, text=mbInk, anchor=west] at (0.60,-5.36)
     {\textbf{24 jetons masqués} — ils étaient cachés, et leurs réponses
      sont notées};

\fill[mbAzure!28, draw=mbAzure!70, line width=0.5pt]
      (0,-5.85) rectangle ++(0.42,-0.42);
\node[font=\sffamily\scriptsize, text=mbInk, anchor=west] at (0.60,-6.06)
     {\textbf{16 jetons de contexte} — ils étaient visibles, et leurs
      réponses sont jetées};
\end{schema}

\begin{intuition}
Imaginez un professeur qui donne un devoir de dix questions. L'élève répond
aux dix. Le professeur, lui, ne corrige que six questions et arrache la page
des quatre autres sans même la lire.

Que va faire l'élève, au bout de quelques devoirs ? Il va comprendre que
quatre questions ne rapportent rien. Il cessera d'y répondre sérieusement.
Il s'en servira peut-être comme brouillon pour préparer les six autres.

Et si un jour on lui demande ce qu'il avait écrit sur ces pages-là, il n'y
aura rien d'utile. Non pas parce qu'il est mauvais, mais parce que personne
ne le lui a jamais demandé.
\end{intuition}

\begin{piege}
Attention à ne pas se tromper de reproche. Ce que nous venons de décrire
n'est pas un bug, ni un oubli des auteurs de V-JEPA~2. C'est un choix, et un
choix qui se défend très bien.

Souvenez-vous du chapitre~\ref{chap:idee-jepa} : le but du jeu est de forcer
l'élève à comprendre la scène pour deviner ce qu'il ne voit pas. Noter les
endroits qu'il voyait déjà semble idiot, puisqu'il lui suffirait de recopier
ce qu'il a sous les yeux. C'est un raisonnement parfaitement logique.

Et pourtant, ce raisonnement a une conséquence que personne n'avait mesurée
avant l'article de V-JEPA~2.1. Nous la découvrirons au chapitre~8, et nous y
mettrons des chiffres. Pour l'instant, contentez-vous de retenir le fait :
\textbf{les jetons de contexte ne reçoivent aucune consigne}.
\end{piege}

\begin{aretenir}
La formule de V-JEPA~2 dit trois choses, et il faut les retenir toutes les
trois :

\begin{enumerate}
  \item on compare des \emph{représentations}, jamais des pixels ;
  \item on mesure l'écart avec la norme L1, c'est-à-dire en effaçant les
        signes et en additionnant ;
  \item et l'on ne regarde que les endroits \emph{masqués}.
\end{enumerate}

Les deux premières sont l'idée JEPA. La troisième est un détail d'écriture
qui tient dans trois symboles, $i \in M$. C'est pourtant elle, et elle
seule, qui va occuper toute la suite de ce livre.
\end{aretenir}

\begin{enia}
Ce petit $i \in M$ n'est pas un détail de mathématicien : c'est une ligne de
programme, et vous l'écrirez vous-même au chapitre~24.

Dans le code, la perte de V-JEPA~2 se calcule en deux temps. On demande
d'abord au prédicteur ses réponses pour tous les jetons, les $2\,048$. Puis
on met de côté celles des endroits masqués, et l'on ne compare que
celles-là. Ce tri, c'est exactement ce que dit $i \in M$.

Deux conséquences, que vous rencontrerez pour de vrai. La première est un
gaspillage : quatre réponses sur dix sont calculées puis jetées, et ce
travail-là, votre machine le paie quand même. La seconde est bien plus
intéressante : \textbf{le premier changement qui mène de V-JEPA~2 à
V-JEPA~2.1 consiste précisément à toucher à ce tri}. On cessera de jeter les
réponses des endroits visibles, et l'on s'en servira au contraire pour
donner une consigne à chaque jeton. Tout le chapitre~10 tient dans cette
seule idée — et il faudra ensuite quelques autres ingrédients, que les
chapitres~11 à~15 vous feront découvrir un par un.
\end{enia}

Il reste une question que nous avons repoussée depuis le début du chapitre.
Cet ensemble $M$, ces trous, ces zones cachées : comment sont-ils choisis ?

\section{Comment on tire le masque au sort}

\begin{enclair}
À chaque pas d'entraînement, on tire un nouveau masque au hasard. Mais pas
n'importe comment : on tire des rectangles, dont on choisit la surface et la
forme dans des plages fixées à l'avance. Et ces rectangles traversent toute
la durée du clip, de la première image à la dernière.
\end{enclair}

Il serait tentant de cacher les jetons un par un, complètement au hasard,
comme on tire des numéros de loto. Ce serait une mauvaise idée, et il faut
comprendre pourquoi.

\begin{intuition}
Si l'on cache un jeton sur deux, en damier, deviner le contenu d'un trou
devient trop facile : il suffit de regarder les quatre voisins, qui sont
tous visibles, et de faire une moyenne. La machine n'a alors besoin de rien
comprendre à la scène : elle apprend seulement à boucher chaque trou avec la
couleur d'à côté.

Pour que le jeu ait un intérêt, il faut cacher des \emph{morceaux entiers},
assez grands pour qu'aucun voisin immédiat ne donne la réponse. La machine
est alors obligée d'aller chercher des indices loin de là, et donc de
comprendre la scène dans son ensemble.
\end{intuition}

Sur ce point, V-JEPA~2.1 n'invente rien : il reprend telle quelle la façon
de faire des modèles V-JEPA qui l'ont précédé. Cette façon de faire repose
sur trois réglages, et nous allons les prendre un par un.

\subsection*{Premier réglage : l'échelle spatiale}

\begin{definition}[L'échelle spatiale du masque][echelle-spatiale]
L'\emph{échelle spatiale} est la part de la surface de l'image que le trou
doit recouvrir. C'est un nombre entre $0$ et $1$ : une échelle de $0{,}3$
veut dire que le trou couvre trente pour cent de la surface.

Dans l'article, cette échelle est tirée au hasard, à chaque pas, entre
$0{,}15$ et $0{,}7$. Autrement dit, le trou recouvre au minimum quinze pour
cent de l'image, et au maximum soixante-dix pour cent.
\end{definition}

\begin{exemple}[De l'échelle à un nombre de cases]
Prenons une image découpée en $16 \times 16 = 256$ cases, ce qui est le
découpage d'une tranche de notre clip habituel.

Avec une échelle de $0{,}15$, le trou doit recouvrir
$0{,}15 \times 256 = 38{,}4$ cases, donc environ trente-huit.

Avec une échelle de $0{,}7$, il doit en recouvrir
$0{,}7 \times 256 = 179{,}2$, donc environ cent soixante-dix-neuf.
\end{exemple}

Nous savons maintenant quelle \emph{surface} le trou doit couvrir. Mais une
même surface peut se ranger de bien des façons : en carré, en rectangle
couché, en rectangle debout. C'est le deuxième réglage qui décide.

\subsection*{Deuxième réglage : le rapport d'aspect}

\begin{definition}[Le rapport d'aspect][rapport-aspect]
Le \emph{rapport d'aspect} d'un rectangle est sa hauteur divisée par sa
largeur. Un rapport de $1$ donne un carré. Un rapport plus petit que $1$
donne un rectangle couché, plus large que haut. Un rapport plus grand que
$1$ donne un rectangle debout, plus haut que large.

Dans l'article, ce rapport est tiré au hasard entre $0{,}75$ et $1{,}5$. Les
trous ne sont donc jamais de longues bandes très fines : ils restent
toujours des rectangles à peu près aussi hauts que larges.
\end{definition}

La figure~\ref{fig:rapport-aspect} montre les trois formes que ce réglage
peut donner, avec les deux valeurs extrêmes que l'article autorise.

\begin{schema}{Ce que le rapport d'aspect change. À gauche, un rapport plus
petit que $1$ : le rectangle est couché, plus large que haut. Au milieu, un
rapport de $1$ : c'est un carré. À droite, un rapport plus grand que $1$ :
le rectangle est debout, plus haut que large. Les trois trous ont pourtant
la même surface : seule leur forme
change.}[rapport-aspect]

% --- couche : rapport 0,75 ---
\fill[mbSlate!35, draw=mbSlate!80, line width=0.7pt] (0,0) rectangle (2.0,1.5);
\draw[<->, draw=mbSlate!70, line width=0.5pt] (0,-0.25) -- (2.0,-0.25);
\node[font=\sffamily\scriptsize, text=mbSlate] at (1.0,-0.55) {largeur};
\draw[<->, draw=mbSlate!70, line width=0.5pt] (-0.25,0) -- (-0.25,1.5);
\node[font=\sffamily\scriptsize, text=mbSlate, rotate=90] at (-0.55,0.75)
     {hauteur};
\node[font=\sffamily\scriptsize\bfseries, text=mbSlate] at (1.0,2.0)
     {rapport $0{,}75$};
\node[font=\sffamily\scriptsize, text=mbAzure!80!black] at (1.0,-1.0)
     {couché};

% --- carre : rapport 1 ---
\begin{scope}[shift={(4.0cm,0cm)}]
  \fill[mbSlate!35, draw=mbSlate!80, line width=0.7pt] (0,0) rectangle (1.73,1.73);
  \node[font=\sffamily\scriptsize\bfseries, text=mbSlate] at (0.87,2.0)
       {rapport $1$};
  \node[font=\sffamily\scriptsize, text=mbAzure!80!black] at (0.87,-1.0)
       {carré};
\end{scope}

% --- debout : rapport 1,5 ---
\begin{scope}[shift={(7.6cm,0cm)}]
  \fill[mbSlate!35, draw=mbSlate!80, line width=0.7pt] (0,0) rectangle (1.41,2.12);
  \node[font=\sffamily\scriptsize\bfseries, text=mbSlate] at (0.7,2.4)
       {rapport $1{,}5$};
  \node[font=\sffamily\scriptsize, text=mbAzure!80!black] at (0.7,-1.0)
       {debout};
\end{scope}
\end{schema}

\begin{exemple}[Lire un rapport d'aspect]
Un trou fait $12$ cases de haut et $16$ cases de large. Son rapport d'aspect
vaut $12 \div 16 = 0{,}75$ : il est plus petit que $1$, donc le trou est
couché. C'est bien ce que disent les deux nombres, puisque $12$ est plus
petit que $16$.

Un autre trou fait $18$ cases de haut et $12$ de large. Son rapport vaut
$18 \div 12 = 1{,}5$ : il est plus grand que $1$, donc le trou est debout.

Et un trou de $14$ cases sur $14$ a pour rapport $14 \div 14 = 1$ : c'est un
carré.
\end{exemple}

Comment passe-t-on de ces deux réglages aux deux côtés du rectangle ? Par
deux petits calculs, toujours les mêmes, et il faut d'abord se rappeler ce
qu'est une racine carrée. Nous l'avions rencontrée au
chapitre~\ref{chap:vit} : la racine carrée d'un nombre, c'est celui qui,
multiplié par lui-même, redonne ce nombre. Ainsi $\sqrt{121} = 11$, puisque
$11 \times 11 = 121$.

Voici les deux calculs. Pour la hauteur, on multiplie la surface par le
rapport, et on prend la racine carrée du résultat. Pour la largeur, on
divise la surface par le rapport, et on prend là aussi la racine carrée.
Ces deux calculs tombent rarement juste : ils donnent des nombres à virgule,
comme $16{,}4$. Or on ne peut pas cacher $16{,}4$ cases : une case est cachée
ou elle ne l'est pas. Il faut donc \emph{arrondir}, c'est-à-dire remplacer le
nombre à virgule par le nombre sans virgule le plus proche. Ainsi $5{,}2$
devient $5$, parce qu'il est plus près de $5$ que de $6$ ; et $7{,}8$ devient
$8$, parce qu'il est plus près de $8$ que de $7$.

\begin{exemple}[Trouver les côtés du rectangle]
Prenons d'abord une surface de $121$ cases avec un rapport de $1$.

\[
  \begin{aligned}
    \text{hauteur} &= \sqrt{121 \times 1} = \sqrt{121} = 11 , \\
    \text{largeur} &= \sqrt{121 \div 1} = \sqrt{121} = 11 .
  \end{aligned}
\]

Le trou fait donc $11$ cases sur $11$. C'est bien un carré, comme le
promettait un rapport de $1$.

Prenons maintenant une surface de $179$ cases avec un rapport de $1{,}5$.

\[
  \begin{aligned}
    \text{hauteur} &= \sqrt{179 \times 1{,}5} = \sqrt{268{,}5} \approx 16{,}4 , \\
    \text{largeur} &= \sqrt{179 \div 1{,}5} = \sqrt{119{,}3} \approx 10{,}9 .
  \end{aligned}
\]

On arrondit : $16{,}4$ devient $16$, et $10{,}9$ devient $11$.

Le trou fait donc $16$ cases de haut sur $11$ de large. Il est bien plus
haut que large, comme le promettait un rapport supérieur à $1$.
\end{exemple}

\begin{piege}
Vérifions la surface obtenue : $16 \times 11 = 176$ cases, alors que nous en
visions $179$. Les deux nombres ne sont pas tout à fait égaux, et c'est
normal : en arrondissant les deux côtés à des nombres entiers, on s'écarte
forcément un peu de la surface visée.

Ne cherchez donc pas à retomber pile sur l'échelle demandée. Elle indique une
taille souhaitée, pas une taille exacte. C'est aussi pourquoi, dans la
figure et dans le programme qui suivent, vous verrez $35$ cases là où
l'échelle $0{,}15$ en demandait $38$ : rien n'est cassé, ce sont seulement
deux arrondis.
\end{piege}

Nos deux premiers réglages dessinent le trou sur \emph{une} image. Reste à
savoir ce qu'il devient quand la vidéo avance. C'est le troisième réglage,
et c'est le plus surprenant des trois.

\subsection*{Troisième réglage : l'échelle temporelle}

\begin{definition}[L'échelle temporelle du masque][echelle-temporelle]
L'\emph{échelle temporelle} est la part de la durée du clip que le trou
recouvre. Dans l'article, elle vaut toujours exactement $1{,}0$ : le trou
traverse donc \emph{toute} la durée du clip, de la première image à la
dernière.
\end{definition}

Ce qui surprend, ici, c'est que ce réglage ne varie jamais. Ce n'est pas un
oubli : c'est un choix, et il est très malin.

\begin{exemple}[Ce que « toujours $1$ » veut dire en nombre d'images]
Un clip de la phase principale contient $16$ images, ce qui fait $8$
tranches de temps une fois les images collées deux par deux.

Une échelle temporelle de $0{,}5$ voudrait dire que le trou ne couvre que
$0{,}5 \times 8 = 4$ tranches, donc la moitié du clip. Avec l'échelle de $1$
retenue par l'article, il en couvre $1 \times 8 = 8$, c'est-à-dire toutes.

Autrement dit, un trou de $150$ cases sur une tranche ne cache pas $150$
jetons, mais $8 \times 150 = 1\,200$ jetons : le même trou est creusé huit
fois, une fois par tranche.
\end{exemple}

\begin{intuition}
Supposons que le trou ne couvre que la moitié du temps. Alors, pour deviner
ce qui est caché à un endroit, la machine n'aurait qu'à regarder le même
endroit un instant plus tôt, quand il était encore visible. Comme les choses
bougent lentement d'une image à l'autre, elle aurait la réponse presque
gratuitement, et sans rien comprendre.

En faisant traverser le trou d'un bout à l'autre du clip, on lui coupe cette
issue de secours. Cet endroit-là n'a \emph{jamais} été visible, à aucun
moment. La machine est donc obligée de raisonner sur ce qui se passe
\emph{autour}, et sur la façon dont la scène évolue. C'est exactement ce
qu'on veut lui apprendre.
\end{intuition}

La figure~\ref{fig:masque-tube} montre cette idée : le trou est une sorte de
tunnel creusé à travers toute la pile d'images.

\begin{schema}{L'échelle temporelle vaut toujours $1$. Les quatre images du
clip sont posées côte à côte, dans l'ordre du temps. Le trou gris tombe
exactement à la même place sur chacune : les traits pointillés le montrent.
Cet endroit de la scène n'est donc visible à aucun moment, et la machine ne
peut pas se contenter de regarder l'image d'avant pour deviner ce qui s'y
cache.}[masque-tube]

% ================== les quatre images, bien séparées ==================
% Chaque image fait 2,1 de large sur 1,6 de haut, et l'on en pose une
% toutes les 2,6 unités. Le trou occupe toujours le même rectangle.
\foreach \k/\nom in {0/{image 1}, 1/{image 2}, 2/{image 3}, 3/{image 4}}{
  \begin{scope}[shift={(\k*2.6,0)}]
    \fill[mbAzure!20, draw=mbAzure!70, line width=0.7pt]
          (0,0) rectangle (2.1,1.6);
    \fill[mbSlate!42, draw=mbSlate!80, line width=0.7pt]
          (0.6,0.45) rectangle (1.5,1.15);
    \node[font=\sffamily\scriptsize, text=mbSlate] at (1.05,-0.32) {\nom};
  \end{scope}
}

% ---- les pointillés : c'est bien le même endroit d'une image à l'autre ----
\draw[draw=mbRose!75, line width=0.7pt, dashed] (0.6,1.15) -- (9.9,1.15);
\draw[draw=mbRose!75, line width=0.7pt, dashed] (0.6,0.45) -- (9.9,0.45);

% ---- la flèche du temps ----
\draw[-{Stealth[round,length=2.6mm,width=2mm]}, draw=mbAmber!85!black,
      line width=1pt] (0,-0.72) -- (9.9,-0.72);
\node[font=\sffamily\scriptsize, text=mbAmber!80!black, anchor=north]
      at (4.95,-0.82) {le temps passe : les images du clip};

% ---- l'annotation, posée au-dessus, sans croiser le dessin ----
\node[font=\sffamily\small\bfseries, text=mbRose!85!black, anchor=west]
      at (0,2.32) {le même trou, à la même place, sur toutes les images};
\node[font=\sffamily\scriptsize, text=mbSlate, anchor=west]
      at (0,1.95) {cet endroit n'est jamais visible, pas une seule fois};
\end{schema}

\begin{surlesdonnees}
Les trois réglages, tels qu'ils sont écrits dans l'annexe de l'article, et
tels que vous les retrouverez au chapitre~24 quand nous coderons :

\begin{center}
\footnotesize
\begin{tabular}{lll}
\toprule
\textbf{Réglage} & \textbf{Valeur} & \textbf{Ce que cela veut dire} \\
\midrule
Échelle spatiale  & de $0{,}15$ à $0{,}7$ &
  le trou couvre de 15 à 70 \% de l'image \\
Échelle temporelle & $1{,}0$ toujours &
  il traverse tout le clip \\
Rapport d'aspect  & de $0{,}75$ à $1{,}5$ &
  ni trop couché, ni trop debout \\
\bottomrule
\end{tabular}
\end{center}

Ces mêmes valeurs sont employées pendant la phase principale et pendant la
phase de refroidissement, sans le moindre changement.
\end{surlesdonnees}

Nous tenons nos trois réglages. Voyons maintenant, de nos yeux, à quoi
ressemblent les trous qu'ils fabriquent.

\subsection*{Trois masques tirés au sort}

Ces trois réglages restent un peu flous tant qu'on n'a pas vu ce qu'ils
donnent. Tirons-en donc trois au hasard, et dessinons les trous obtenus.

La figure~\ref{fig:trois-masques} montre le résultat sur une tranche de
notre clip habituel, c'est-à-dire sur une grille de $16 \times 16$ cases.
Les trois tirages sont ceux que nous ferons calculer par le programme, tout
à l'heure : vous retrouverez donc exactement les mêmes nombres.

\begin{schema}{Trois masques tirés au sort, dessinés sur une tranche de
$16 \times 16$ cases. À gauche, le plus petit trou que l'article autorise ;
au milieu, un trou de taille moyenne, ici carré ; à droite, le plus gros,
un rectangle debout si haut qu'il touche les deux bords de la grille. Les
trois viennent des mêmes deux réglages, tirés au hasard dans leurs plages :
seuls les nombres tirés changent.}[trois-masques]

% ================= panneau 1 : petit trou, 5 x 7 =================
\begin{scope}[shift={(0cm,0cm)}]
  \begin{scope}[x=0.15cm, y=0.15cm]
    \fill[mbAzure!18] (0,0) rectangle (16,-16);
    \foreach \k in {1,...,15}{
      \draw[draw=mbAzure!40, line width=0.2pt] (\k,0) -- (\k,-16);
      \draw[draw=mbAzure!40, line width=0.2pt] (0,-\k) -- (16,-\k);
    }
    \fill[mbSlate!45] (5,-4) rectangle (12,-9);
    \draw[draw=mbSlate!85, line width=0.6pt, dashed] (5,-4) rectangle (12,-9);
    \draw[draw=mbAzure!70, line width=0.7pt] (0,0) rectangle (16,-16);
  \end{scope}
  \node[font=\sffamily\scriptsize\bfseries, text=mbSlate] at (1.2,-2.85)
       {petit trou};
  \node[font=\sffamily\scriptsize, text=mbSlate, align=center] at (1.2,-3.5)
       {échelle $0{,}15$\\rapport $0{,}75$};
  \node[font=\sffamily\scriptsize, text=mbAzure!80!black] at (1.2,-4.2)
       {$5 \times 7 = 35$ cases};
\end{scope}

% ================= panneau 2 : trou moyen, 11 x 11 =================
\begin{scope}[shift={(3.3cm,0cm)}]
  \begin{scope}[x=0.15cm, y=0.15cm]
    \fill[mbAzure!18] (0,0) rectangle (16,-16);
    \foreach \k in {1,...,15}{
      \draw[draw=mbAzure!40, line width=0.2pt] (\k,0) -- (\k,-16);
      \draw[draw=mbAzure!40, line width=0.2pt] (0,-\k) -- (16,-\k);
    }
    \fill[mbSlate!45] (2,-2) rectangle (13,-13);
    \draw[draw=mbSlate!85, line width=0.6pt, dashed] (2,-2) rectangle (13,-13);
    \draw[draw=mbAzure!70, line width=0.7pt] (0,0) rectangle (16,-16);
  \end{scope}
  \node[font=\sffamily\scriptsize\bfseries, text=mbSlate] at (1.2,-2.85)
       {trou moyen};
  \node[font=\sffamily\scriptsize, text=mbSlate, align=center] at (1.2,-3.5)
       {échelle $0{,}47$\\rapport $1{,}00$};
  \node[font=\sffamily\scriptsize, text=mbAzure!80!black] at (1.2,-4.2)
       {$11 \times 11 = 121$ cases};
\end{scope}

% ================= panneau 3 : gros trou, 16 x 11 =================
\begin{scope}[shift={(6.6cm,0cm)}]
  \begin{scope}[x=0.15cm, y=0.15cm]
    \fill[mbAzure!18] (0,0) rectangle (16,-16);
    \foreach \k in {1,...,15}{
      \draw[draw=mbAzure!40, line width=0.2pt] (\k,0) -- (\k,-16);
      \draw[draw=mbAzure!40, line width=0.2pt] (0,-\k) -- (16,-\k);
    }
    \fill[mbSlate!45] (3,0) rectangle (14,-16);
    \draw[draw=mbSlate!85, line width=0.6pt, dashed] (3,0) rectangle (14,-16);
    \draw[draw=mbAzure!70, line width=0.7pt] (0,0) rectangle (16,-16);
  \end{scope}
  \node[font=\sffamily\scriptsize\bfseries, text=mbSlate] at (1.2,-2.85)
       {gros trou};
  \node[font=\sffamily\scriptsize, text=mbSlate, align=center] at (1.2,-3.5)
       {échelle $0{,}70$\\rapport $1{,}50$};
  \node[font=\sffamily\scriptsize, text=mbAzure!80!black] at (1.2,-4.2)
       {$16 \times 11 = 176$ cases};
\end{scope}
\end{schema}

\begin{exemple}[Les trois masques, en jetons]
Chacun de ces trois trous traverse les $8$ tranches de temps du clip,
puisque l'échelle temporelle vaut toujours $1$. Il suffit donc de
multiplier le nombre de cases par $8$ pour obtenir le nombre de jetons
cachés :

\begin{center}
\footnotesize
\begin{tabular}{lrrr}
\toprule
\textbf{Tirage} & \textbf{Cases} & \textbf{Jetons masqués} &
\textbf{Jetons de contexte} \\
\midrule
petit trou  & $35$  & $280$   & $1\,768$ \\
trou moyen  & $121$ & $968$   & $1\,080$ \\
gros trou   & $176$ & $1\,408$ & $640$ \\
\bottomrule
\end{tabular}
\end{center}

D'un tirage à l'autre, le nombre de jetons cachés passe donc de $280$ à
$1\,408$, soit cinq fois plus. La difficulté du jeu change beaucoup d'un pas
d'entraînement au suivant, et c'est voulu : la machine doit savoir se
débrouiller aussi bien avec quelques indices qu'avec beaucoup.
\end{exemple}

\begin{surlesdonnees}
Regardez bien le panneau de droite. Le calcul demandait une hauteur de
$16{,}4$ cases, qui une fois arrondie donne $16$ : c'est exactement la
hauteur de la grille. Le trou touche donc le bord du haut et le bord du bas.

Il s'en est fallu de peu que le rectangle ne dépasse. C'est pour ce cas-là
que notre programme contiendra un \co{python}{min} : si jamais le calcul
réclame plus de cases que la grille n'en possède, on ramène simplement le
rectangle à la taille de la grille. Ce n'est ni une erreur ni un cas rare,
et le programme continue sans se plaindre.
\end{surlesdonnees}

\begin{remarque}
Dans la pratique, on ne se contente pas d'un seul rectangle. On en tire
plusieurs, et l'on cache tout ce qui est recouvert par au moins l'un d'eux.
Quand deux rectangles se chevauchent, la partie commune n'est évidemment
comptée qu'une seule fois : c'est ce qu'on appelle prendre leur
\emph{réunion}. Le résultat est un trou de forme moins régulière, ce qui
rend le jeu plus intéressant. Nous le ferons tourner dans la section de code,
et vous verrez apparaître les nombres qui nous accompagnent depuis le
chapitre~\ref{chap:idee-jepa}.
\end{remarque}

Nous avons la formule et nous avons le masque. Vérifions tout cela sur une
machine.

\section{Passons au code}

Nous allons écrire quatre petits programmes. Le premier refait, ligne par
ligne, le calcul de perte que nous venons de mener à la main : vous devez
retrouver le $0{,}6$, le $0{,}4$ et le $0{,}5$. Le deuxième compare la norme
L1 aux carrés du chapitre~\ref{chap:apprendre}. Les deux derniers s'occupent
du masque : l'un calcule la taille d'un trou à partir de ses réglages,
l'autre pose deux trous sur une grille et compte les jetons cachés.

Commençons par la perte.

\begin{code}{python}[La perte de V-JEPA~2, sur nos deux jetons]
import torch

# Deux jetons masques, de trois nombres chacun.
# Ligne 0 : le jeton numero 3. Ligne 1 : le jeton numero 7.
prediction = torch.tensor([[0.5, -0.2, 1.0],
                           [0.0,  0.0, 0.0]])
cible      = torch.tensor([[0.4,  0.1, 0.8],
                           [0.1, -0.2, 0.1]])

# 1. l'ecart, terme a terme, puis on efface les signes
ecarts = (prediction - cible).abs()
print(ecarts)

# 2. la norme L1 : on additionne les nombres de chaque jeton
par_jeton = ecarts.sum(dim=1)
print(par_jeton)

# 3. la moyenne sur les jetons masques : c'est la perte
perte = par_jeton.mean()
print(perte)
\end{code}

Voici ce que le programme affiche :

\begin{sortie}
tensor([[0.1000, 0.3000, 0.2000],
        [0.1000, 0.2000, 0.1000]])
tensor([0.6000, 0.4000])
tensor(0.5000)
\end{sortie}

Reprenez le tableau de la section~\ref{sec:calcul-complet} et comparez. Le premier affichage donne
les écarts, une fois les signes effacés, exactement comme nous les avions
calculés. Le deuxième donne les deux normes L1, soit $0{,}6$ puis $0{,}4$.
Le troisième donne la perte, $0{,}5$. La machine n'a rien fait de plus que
ce que nous venons de faire au crayon.

\begin{piege}
Regardez bien la ligne \co{python}{ecarts.sum(dim=1)}. Le réglage
\co{python}{dim=1} dit : « additionne le long du deuxième axe », c'est-à-dire
à l'intérieur de chaque jeton. Si vous écriviez \co{python}{dim=0}, vous
additionneriez d'un jeton à l'autre, ce qui n'a aucun sens ici. Le
programme ne se plaindrait pas : il rendrait trois nombres au lieu de deux,
et la perte serait fausse en silence. Reprenez le réflexe du
chapitre~\ref{chap:image-video} : affichez la forme, et demandez-vous si
elle est celle que vous attendiez.
\end{piege}

Passons à la comparaison entre la norme L1 et les carrés. C'est le calcul de
l'encadré « pour aller plus loin », vérifié sur machine.

\begin{code}{python}[Norme L1 contre carrés : qui punit le plus les gros écarts ?]
import torch

ecarts = torch.tensor([0.1, 1.0])     # un petit ecart, et un gros

total_l1     = ecarts.abs().sum()
total_carres = (ecarts ** 2).sum()

print(f"norme L1 : total = {total_l1:.2f}")
print(f"   le petit ecart compte pour {100 * 0.1 / total_l1:.0f} %")
print(f"   le gros  ecart compte pour {100 * 1.0 / total_l1:.0f} %")

print(f"carres   : total = {total_carres:.2f}")
print(f"   le petit ecart compte pour {100 * 0.01 / total_carres:.0f} %")
print(f"   le gros  ecart compte pour {100 * 1.00 / total_carres:.0f} %")
\end{code}

Et le résultat :

\begin{sortie}
norme L1 : total = 1.10
   le petit ecart compte pour 9 %
   le gros  ecart compte pour 91 %
carres   : total = 1.01
   le petit ecart compte pour 1 %
   le gros  ecart compte pour 99 %
\end{sortie}

Avec les carrés, le petit écart ne pèse plus qu'un pour cent : autant dire
qu'il a disparu. Avec la norme L1, il pèse encore neuf fois plus. C'est
précisément ce que V-JEPA~2 cherche : continuer d'écouter les petits écarts,
au lieu de ne s'occuper que des grosses surprises imprévisibles.

Terminons par le masque. Nous allons d'abord calculer la taille d'un
rectangle à partir de ses deux réglages, puis en poser deux sur une grille
et compter.

\begin{code}{python}[Calculer la taille d'un trou à partir de ses réglages]
def taille_du_trou(echelle, rapport, H=16, W=16):
    """Rend la hauteur et la largeur du rectangle, en cases."""
    aire = echelle * H * W                       # la surface visee
    hauteur = round((aire * rapport) ** 0.5)
    largeur = round((aire / rapport) ** 0.5)
    return min(hauteur, H), min(largeur, W)      # on ne sort pas de la grille

for echelle, rapport in [(0.15, 0.75), (0.473, 1.0), (0.70, 1.5)]:
    h, w = taille_du_trou(echelle, rapport)
    print(f"echelle {echelle:.3f}, rapport {rapport:.2f} -> "
          f"{h:2d} x {w:2d} = {h * w:3d} cases, soit {8 * h * w:4d} jetons")
\end{code}

Le programme affiche :

\begin{sortie}[Trois trous, du plus petit au plus grand]
echelle 0.150, rapport 0.75 ->  5 x  7 =  35 cases, soit  280 jetons
echelle 0.473, rapport 1.00 -> 11 x 11 = 121 cases, soit  968 jetons
echelle 0.700, rapport 1.50 -> 16 x 11 = 176 cases, soit 1408 jetons
\end{sortie}

Le facteur $8$ de la dernière colonne, c'est le nombre de tranches de temps
d'un clip de la phase principale, que nous avions calculé au
chapitre~\ref{chap:vit} : $16$ images divisées par des tubelets de $2$ font
bien $8$ tranches. Comme le trou traverse tout le clip, on multiplie le
nombre de cases par ce $8$.

Remarquez la dernière ligne : la hauteur voulue dépassait la grille, et le
\co{python}{min} l'a ramenée à $16$. C'est normal, et le programme continue
sans se plaindre.

Posons enfin deux rectangles sur la grille, et comptons.

\begin{code}{python}[Fabriquer un masque et compter les jetons cachés]
import torch

T, H, W = 8, 16, 16          # tranches de temps, lignes, colonnes

masque = torch.zeros(H, W, dtype=torch.bool)   # au depart, rien n'est cache
masque[2:13, 2:13] = True                      # premier bloc : 11 x 11
masque[9:16, 10:16] = True                     # second bloc  :  7 x  6

par_tranche = int(masque.sum())      # la reunion des deux blocs
total = par_tranche * T              # le trou traverse tout le clip

print(f"cases cachees sur une tranche : {par_tranche}")
print(f"jetons masques  (|M|)         : {total}")
print(f"jetons visibles (le contexte) : {T * H * W - total}")
print(f"part jetee par la perte       : "
      f"{100 * (T * H * W - total) / (T * H * W):.1f} %")
\end{code}

Le résultat :

\begin{sortie}[Le compte final]
cases cachees sur une tranche : 151
jetons masques  (|M|)         : 1208
jetons visibles (le contexte) : 840
part jetee par la perte       : 41.0 %
\end{sortie}

Arrêtons-nous sur ces quatre lignes, car elles bouclent le chapitre.

Le premier bloc couvre $11 \times 11 = 121$ cases, le second $7 \times 6 =
42$. Cela ferait $163$ si on les additionnait bêtement. Mais ils se
chevauchent sur $12$ cases, et ces cases-là ne doivent être comptées qu'une
fois : le total est donc $121 + 42 - 12 = 151$. C'est bien ce qu'affiche le
programme, et c'est ce que veut dire « prendre la réunion ».

Ces $151$ cases, multipliées par les $8$ tranches de temps, donnent $1\,208$
jetons masqués et $840$ jetons de contexte. Ce sont, à quelques unités près,
les $1\,200$ et $848$ qui nous accompagnent depuis le
chapitre~\ref{chap:idee-jepa} — et nous continuerons d'employer ces
chiffres ronds, plus faciles à retenir.

Et la dernière ligne dit tout : \textbf{quarante et un pour cent des réponses
du prédicteur sont fabriquées, puis jetées}.

\begin{aretenir}
L'équation~1 de V-JEPA~2, en une phrase de français : \emph{on mesure, aux
seuls endroits masqués, l'écart entre ce que le prédicteur a répondu et ce
que le professeur avait produit, on efface les signes, on additionne, et on
divise par le nombre d'endroits masqués}.

Et le mot qui compte dans cette phrase est le mot « seuls ».
\end{aretenir}

\vspace{0.5em}

Nous voici au bout de la partie~II. Nous savons comment la machine est
bâtie, pourquoi elle ne s'écroule pas, et ce qu'elle cherche exactement à
faire baisser. Toutes les pièces de V-JEPA~2 sont maintenant entre vos
mains.

Cette formule marche d'ailleurs très bien pour reconnaître une action dans
une vidéo, et l'article de V-JEPA~2 l'a prouvé sur des dizaines de mesures.
Mais que se passe-t-il si on lui demande autre chose ? Si, au lieu de
demander « que se passe-t-il dans cette vidéo ? », on demande « à quelle
distance se trouve ce pixel-là ? », ou bien « à quel objet appartient ce
pixel-là ? »

Ces questions-là ne réclament pas une seule réponse pour toute l'image :
elles en réclament une \emph{par endroit}. Pour y répondre, il faudrait que
chaque jeton décrive bien sa propre place — y compris les jetons de
contexte, ceux dont personne, nous venons de le voir, ne demande jamais
rien.

Vous devinez déjà où cela nous mène : sur ces questions-là, la formule que
nous venons d'écrire échoue. Et c'est précisément cet échec qui a donné
naissance à V-JEPA~2.1. Avant de le constater nous-mêmes, il faut apprendre
à poser ces nouvelles questions et à noter les réponses : c'est le sujet du
chapitre suivant, qui ouvre la partie~III.

\phantomsection
\section*{Exercices}
\addcontentsline{toc}{section}{Exercices}

\begin{exercice}[Calculer une norme L1]
Un prédicteur répond $(1{,}2 \ ; \ -0{,}5 \ ; \ 0{,}3 \ ; \ 0{,}0)$ alors
que la cible valait $(1{,}0 \ ; \ -0{,}1 \ ; \ 0{,}3 \ ; \ -0{,}2)$. Quelle
est la norme L1 de l'écart ?
\end{exercice}

\begin{solution}
On soustrait terme à terme :
\[
  1{,}2 - 1{,}0 = 0{,}2 , \quad
  -0{,}5 + 0{,}1 = -0{,}4 , \quad
  0{,}3 - 0{,}3 = 0 , \quad
  0{,}0 + 0{,}2 = 0{,}2 .
\]

On efface les signes, ce qui donne $(0{,}2 \ ; \ 0{,}4 \ ; \ 0 \ ; \
0{,}2)$, puis on additionne :
\[
  0{,}2 + 0{,}4 + 0 + 0{,}2 = 0{,}8 .
\]

La norme L1 vaut donc $0{,}8$. Remarquez le troisième nombre : la machine
était tombée parfaitement juste dessus, et il ne coûte donc rien du tout.
\end{solution}

\begin{exercice}[Une perte complète sur trois jetons]
Le masque recouvre trois jetons. Les normes L1 de leurs écarts valent
$1{,}2$, $0{,}9$ et $0{,}3$. Que vaut la perte ? Et si un quatrième jeton
était masqué, avec un écart de $0{,}0$, que deviendrait-elle ?
\end{exercice}

\begin{solution}
Avec trois jetons, $|M| = 3$ et la somme vaut $1{,}2 + 0{,}9 + 0{,}3 =
2{,}4$. La perte vaut donc
\[
  \frac{2{,}4}{3} = 0{,}8 .
\]

Avec un quatrième jeton d'écart nul, la somme ne bouge pas — elle vaut
toujours $2{,}4$ — mais $|M|$ passe à $4$. La perte devient
\[
  \frac{2{,}4}{4} = 0{,}6 .
\]

Elle a baissé, alors que la machine ne s'est pas améliorée d'un pouce sur
les trois premiers jetons. C'est l'effet de la division : la perte est une
\emph{moyenne}, et ajouter une bonne réponse tire la moyenne vers le bas.
C'est aussi pourquoi on ne compare jamais deux pertes calculées sur des
masques de tailles différentes.
\end{solution}

\begin{exercice}[Compter les réponses jetées]
Un clip de la phase de refroidissement contient $18\,432$ jetons. Le masque
en recouvre $11\,000$. Combien de réponses le prédicteur produit-il ?
Combien la perte en regarde-t-elle ? Combien sont jetées, et quelle part
cela représente-t-il ?
\end{exercice}

\begin{solution}
Le prédicteur produit une réponse par jeton reçu, et il reçoit le contexte
plus les masques, c'est-à-dire tout le clip : il produit donc $18\,432$
réponses.

La perte ne regarde que les jetons masqués, soit $11\,000$ réponses.

Il en reste donc $18\,432 - 11\,000 = 7\,432$ qui sont jetées, ce qui
représente
\[
  \frac{7\,432}{18\,432} \approx 0{,}403 ,
\]
soit un peu plus de quarante pour cent. On retrouve, à très peu de chose
près, la même proportion que sur un clip de la phase principale : ce n'est
pas la taille du clip qui décide de cette part, mais la taille du masque.
\end{solution}

\begin{exercice}[Retrouver la taille d'un trou]
Sur une grille de $16 \times 16$ cases, on tire une échelle spatiale de
$0{,}25$ et un rapport d'aspect de $1{,}0$. Combien de cases le trou
couvre-t-il, et quelles sont ses dimensions ? Et comme le clip compte $8$
tranches de temps, combien de jetons cela fait-il en tout ?
\end{exercice}

\begin{solution}
La surface visée vaut $0{,}25 \times 256 = 64$ cases. On applique ensuite
les deux calculs habituels. Hauteur : $\sqrt{64 \times 1} = 8$. Largeur :
$\sqrt{64 \div 1} = 8$. Le trou fait donc $8$ cases sur $8$, et c'est bien
un carré, comme le voulait un rapport de $1$.

Comme l'échelle temporelle vaut toujours $1$, ce même carré est creusé sur
les huit tranches de temps. Cela fait
\[
  8 \times 64 = 512 \text{ jetons masqués,}
\]
et il reste donc $2\,048 - 512 = 1\,536$ jetons de contexte.
\end{solution}

\begin{exercice}[Pourquoi le trou traverse tout le clip]
Un collègue propose de ne masquer que la moitié des images du clip, pour
« laisser un peu d'aide à la machine ». Pourquoi est-ce une mauvaise idée ?
\end{exercice}

\begin{solution}
Parce que cette « aide » est bien trop grande, et qu'elle vide le jeu de son
intérêt.

Si un endroit est visible sur la première moitié du clip et caché sur la
seconde, la machine n'a plus besoin de comprendre quoi que ce soit : il lui
suffit de recopier ce qu'elle a vu au même endroit un instant plus tôt.
Comme les choses bougent peu d'une image à l'autre, elle aura presque juste
à tous les coups, et la perte descendra très bas sans qu'elle ait rien
appris.

C'est exactement pour fermer cette porte que l'article fixe l'échelle
temporelle à $1{,}0$ et ne la fait jamais varier. L'endroit caché ne doit
avoir été visible à aucun moment.
\end{solution}

\begin{exercice}[Le mot qui compte]
Résumez en une phrase ce que la formule de V-JEPA~2 fait des réponses que le
prédicteur produit pour les jetons de contexte. Puis dites, en une phrase,
pourquoi ce choix pouvait sembler raisonnable à ses auteurs.
\end{exercice}

\begin{solution}
Elle les jette : le signe $\sum_{i \in M}$ ne fait prendre à $i$ que les
numéros des jetons masqués, si bien que les réponses produites aux endroits
visibles ne sont jamais comparées à quoi que ce soit, et ne comptent donc
pour rien dans la perte.

Ce choix semblait raisonnable parce que noter la machine là où elle voit
déjà tout paraît sans intérêt : il lui suffirait de recopier son entrée pour
avoir juste, et l'on n'apprendrait rien. Le raisonnement est logique — mais
nous verrons au chapitre~8 qu'il a une conséquence que personne n'avait
mesurée, et au chapitre~10 comment V-JEPA~2.1 la corrige.
\end{solution}
